% Related Work: Decentralized Optimization and Matrix‑Aware Optimizers
% This file is auto‑generated. Edit with care.

\section{Related Work}
\label{sec:related-work}

In this section we review (i) matrix–aware optimizers with an emphasis on Muon and its variants, (ii) decentralized stochastic optimization and heterogeneity–correction mechanisms, and (iii) recent attempts at combining matrix orthogonalization with decentralized training. This discussion is organized to highlight the algorithmic and theoretical gaps that motivate our SUDA–Muon framework.

\subsection{Muon and Matrix–Aware Optimizers}

Classical first–order methods such as SGD and AdamW treat all trainable weights as vectors equipped with a Euclidean geometry, ignoring that most parameters in modern architectures are structured matrices (e.g., attention projections, MLP kernels). This mismatch has motivated a growing line of work on matrix–aware or geometry–aware optimizers that operate under operator‐norm or manifold constraints rather than coordinatewise scaling.

\paragraph{Matrix orthogonalization and Muon.}
The Muon optimizer, originally introduced as a steepest–descent method under spectral–norm constraints, replaces the raw gradient $G$ with its polarized direction $\operatorname{msgn}(G)$ obtained from the matrix sign or, equivalently, the unitary factor of the polar decomposition. This orthogonalization makes the update direction scale–invariant and better aligned with the operator geometry of each layer. Subsequent work has refined both the numerical and theoretical aspects of this idea. Amsel et~al.\ provide optimal matrix–sign solvers tailored to Muon, showing how Chebyshev–type polynomial accelerations and low–degree factorizations reduce the cost of Newton–Schulz iterations while preserving stability of the orthogonal factor. Sato et~al. analyze nonconvex convergence and critical batch size of Muon, connecting its dynamics to classical large–batch scaling laws and demonstrating that orthogonalization can delay the onset of the generalization–efficiency tradeoff. Kovalev studies gradient orthogonalization through a non–Euclidean trust–region lens, establishing global convergence guarantees for trust–region methods whose local model is built in orthogonalized directions, and clarifying when orthogonalization acts as an implicit preconditioner versus a constraint enforcement mechanism.

On the empirical side, multiple works have confirmed that Muon scales competitively in large–scale language modeling. Liu et~al. show that Muon can pretrain billion–parameter LLMs with improved data efficiency when combined with carefully tuned weight decay and learning–rate schedules. Shah et~al. conduct a systematic comparison of Muon and AdamW across model sizes and batch regimes, finding that Muon consistently expands the Pareto frontier in terms of tokens~vs.~loss and remains robust at batch sizes beyond the ``critical batch size'' for AdamW. Chen et~al. interpret Muon as optimization under spectral–norm constraints, providing a principled perspective on why orthogonalization mitigates exploding activations and improves hyperparameter transfer across widths. More broadly, Bernstein and Newhouse, as well as Yang et~al., place Muon–style methods within a family of non–Euclidean steepest–descent algorithms governed by layerwise norms and spectral conditions, helping to formalize the design space of matrix–aware optimizers.

The basic Muon update has also been extended in several directions. AdaMuon and other adaptive variants inject element–wise second–moment estimates on top of the orthogonalized direction, seeking to blend curvature–aware scaling with operator–norm constraints. NorMuon, LiMuon, and MuonBP focus on reducing the computational overhead of exact orthogonalization via block–periodic schemes, low–rank projections, or dimension–reduced matrix signs, while retaining most of the stability benefits. TEON and related tensorized approaches generalize layerwise orthogonalization to structured tensor factorizations, enabling Muon–style updates on factored or low–rank parameterizations. PolarGrad further unifies Muon, Scion, and other matrix–gradient optimizers as instances of preconditioned steepest descent in which the preconditioner encodes a polar–decomposition–like normalization.

Another branch of work explores Muon beyond centralized LLM pretraining. FedMuon and related studies investigate matrix orthogonalization in federated learning, arguing that scale–invariant local updates can mitigate client–drift under heterogeneous data. Dion and its successors (CurvaDion, Dion2) design distributed orthonormalized updates compatible with sharded model partitions, largely targeting data–parallel LLM training on datacenter clusters. These methods, however, typically assume synchronous all–reduce style communication and retain an effectively centralized optimization objective; they do not address the algorithmic challenges posed by fully decentralized, graph–constrained optimization or by consensus constraints on the model parameters themselves.

\paragraph{Limitations of existing matrix–aware optimizers.}
Despite this rapid progress, current matrix–aware optimizers are essentially single–node or parameter–server methods. Their convergence analyses assume a centralized objective and homogeneous noise, and their communication models rely on dense, global synchronization. The nonlinearity of the matrix sign (or more general orthogonalization maps) also makes it nontrivial to embed Muon into architectures where the optimization variables are implicitly constrained through consensus or mixing operators, as in decentralized networks. In particular, existing theory does not characterize how orthogonalized updates interact with the bias–correction mechanisms and topology–dependent error terms that arise in decentralized stochastic gradient methods.

\subsection{Decentralized Optimization and Heterogeneity Correction}

Decentralized parallel SGD (DPSGD) and related methods have emerged as a powerful alternative to parameter–server architectures, alleviating central bottlenecks and improving robustness by allowing each agent to communicate only with its neighbors. In the basic setting, $n$ agents seek to minimize the average of local objectives $f_i$ using local gradients and a mixing matrix $W$ that encodes the communication graph. However, when data across agents are non–IID, local objectives can differ significantly, leading to the well–known phenomenon of client drift: local iterates move towards different minimizers and consensus is lost.

\paragraph{Bias correction via EXTRA and Exact Diffusion.}
To correct this drift, a rich line of work has developed bias–correction schemes that modify the raw consensus updates using carefully designed combinations of current and past iterates. The EXTRA algorithm of Shi et~al. introduces an exact first–order method that achieves linear convergence under constant step sizes by replacing the standard consensus step with a two–step recursion involving both $W$ and $W^2$. This effectively cancels the steady–state bias of vanilla DPSGD and yields exact convergence to the network optimum. Subsequent extensions provide variance–reduced and accelerated variants of EXTRA and DIGing, and extend these ideas to composite optimization and proximal settings.

The Exact Diffusion (ED) or $D^2$ framework developed by Alghunaim, Yuan, Sayed and collaborators refines this perspective. ED reinterprets bias correction as a diffusion process on the network, with carefully chosen correction terms that decouple the impact of data heterogeneity from the mixing topology. A unified and refined convergence analysis for nonconvex decentralized learning shows that, for a broad class of algorithms including ED and decentralized SGD, the dominant convergence rate terms can be decomposed into data–dependent and topology–dependent components, where the latter enter only through higher–order powers of the mixing matrix. This separation underlies network–independent step–size rules and enables linear speedup results in which the effective sample complexity scales nearly linearly with the number of agents.

\paragraph{Gradient tracking and primal–dual unification.}
Gradient tracking (GT) methods pursue a complementary strategy in which each agent maintains an auxiliary variable tracking the global gradient. Early work by Pu and Nedi\'c and many subsequent developments show that GT can achieve linear convergence for strongly convex objectives and sublinear rates in the nonconvex setting, even under communication compression, local updates, or time–varying graphs. Recent analyses quantify the trade–offs between computation and communication, study optimal parameter choices, and extend GT to manifold–constrained problems and min–max games.

More recently, primal–dual viewpoints and unified frameworks have been proposed. Alghunaim and Yuan's refined analysis can be interpreted as a primal–dual recursion driven by low–degree polynomials of the mixing matrix. Frameworks such as SPARKLE and other single–loop primal–dual methods further demonstrate that gradient tracking, EXTRA, ED, and related schemes can all be embedded into a common template parameterized by a handful of polynomial and step–size coefficients. These unifying perspectives clarify how different algorithmic choices trade bias correction for robustness to heterogeneity and how network topology affects second–order error terms but not the leading–order convergence rate.

Despite this progress, the vast majority of decentralized methods are built around Euclidean, element–wise gradient updates. Even when manifold constraints are considered (e.g., decentralized optimization on the Stiefel manifold or other compact submanifolds), the emphasis is on Riemannian gradients and retractions rather than on nonlinear matrix orthogonalization mappings akin to the matrix sign. As a result, existing theories and algorithm designs do not directly accommodate the type of structured, non–Euclidean updates required by Muon.

\subsection{Decentralized Muon and the Need for a Unified Framework}

Only very recently have researchers begun to explore decentralized variants of Muon. DeMuon extends Muon to decentralized graphs by combining local Newton–Schulz orthogonalization with a gradient–tracking consensus mechanism. Each agent applies matrix polarization to its local gradient and then participates in a GT–style correction step to mitigate heterogeneity, thereby approximating centralized Muon on the network average objective. Related works such as Dion address distributed orthonormalized updates, but in a data–parallel rather than consensus–constrained setting, typically assuming an all–reduce communication primitive.

While these initial attempts demonstrate that orthogonalized updates can be implemented over graphs, they suffer from several limitations that motivate our SUDA–Muon framework.

\paragraph{Algorithmic rigidity.}
Existing decentralized Muon designs are tied to a single communication template, most notably gradient tracking. This rigidity makes it difficult to exploit other powerful heterogeneity–correction mechanisms such as Exact Diffusion, EXTRA, or more general primal–dual recursions, and it obscures the role of the mixing matrix polynomial that implicitly governs error propagation. In contrast, SUDA–style frameworks show that many decentralized Euclidean algorithms can be viewed as instantiations of a low–degree polynomial in the mixing operator applied to suitably defined primal–dual variables. Our work extends this unifying perspective to the non–Euclidean setting where local updates are produced by matrix polarization rather than standard gradients.

\paragraph{Lack of topology–separated convergence analysis.}
Current analyses of decentralized Muon algorithms largely inherit proof techniques from classical GT methods without exploiting the refined topology–separated convergence phenomena identified in SUDA and related unified analyses. As a consequence, they typically yield conservative step–size conditions and convergence rates that entangle data heterogeneity, noise, and spectral properties of the communication graph. By formulating SUDA–Muon as a primal–dual recursion driven by low–degree polynomials of the mixing matrix, we recover convergence bounds in which network topology appears only in higher–order error terms, mirroring the behavior of Euclidean SUDA while accounting for the nonlinear matrix orthogonalization map.

\paragraph{Open questions on linear speedup.}
A central property of decentralized SGD is its asymptotic linear speedup: under suitable conditions, the time to reach a target accuracy decreases nearly linearly with the number of agents. For Muon, however, the interplay between nonlinear polarization and network averaging remains poorly understood. Existing decentralized Muon schemes do not establish whether linear speedup is achievable, nor do they clarify how the spectral norm constraints induced by matrix orthogonalization affect noise averaging and variance reduction across the network. Our analysis shows that, under mild assumptions on the polarization operator and mixing matrix, SUDA–Muon inherits the linear–speedup guarantees of its Euclidean counterparts while maintaining the scale–invariance and stability benefits of Muon.

Taken together, the literature on matrix–aware optimizers and decentralized optimization suggests that these two lines of work have so far evolved largely in parallel. Muon and its variants deliver strong performance in centralized or parameter–server settings by respecting matrix geometry, while SUDA–style frameworks offer a unifying and topology–robust view of decentralized Euclidean optimization. The goal of SUDA–Muon is to bridge this gap: to design a family of decentralized algorithms that simultaneously (i) exploit matrix orthogonalization at each agent, (ii) fall within a unified primal–dual template encompassing EXTRA, Exact Diffusion, and gradient tracking, and (iii) admit refined, topology–separated convergence guarantees including linear speedup with respect to the number of agents.
